This dissertation is presented in a cumulative format, with the topics necessary to understand the appended publications briefly summarized in the following sections. \cref{sec:nmr} introduces the fundamentals of \gls{MRI} signal generation and image encoding, which are subsequently extended in \cref{sec:spiralacq} to encompass non-Cartesian encoding and image reconstruction. \cref{sec:bSSFP} provides a brief overview of the taxonomy of rapid steady-state sequences and the signal characteristics of \gls{bSSFP} contrast. The two neuroscience applications addressed in \hyperref[sec:paper1]{Publication 1} and \hyperref[sec:paper2]{2}, namely \gls{fMRI} and \gls{DMI}, are subsequently presented in \cref{sec:fmri} and \cref{sec:DMI}, respectively. The research objectives and principal findings are summarized in the forthcoming sections, and the dissertation concludes with an outlook that considers the potential and prospective directions of the developed methodology.

\section{NMR Signal} \label{sec:nmr}
The \gls{NMR} signal arises from the intrinsic angular momentum of the atomic nuclei also known as spin, which is a quantum mechanical property. Individual nuclei with non-zero spin are associated with a magnetic moment due to their charge. However, these magnetic moments are aligned randomly  in the absence of an external magnetic field, resulting in no net magnetic moment. According to quantum mechanics, when the spins experience an external magnetic field $B_0$, the energy state of the spin system undergoes splitting known as the \emph{Zeeman effect} into $2s+1$ discrete energy levels, where $s$ is the spin quantum number of the nucleus. Without other interactions, these energy levels are $\Delta E=\hbar\gamma B_0$ apart~\cite{levittSpinDynamicsBasics}, where $\gamma$ is the gyromagnetic ratio for the nucleus and $\hbar$ is the reduced Planck's constant. The energy difference also implies that the spins are precessing around an external magnetic field at the \emph{Larmor frequency} $\omega_0= \gamma B_0$. The electronic environment near the nuclei can influence the magnetic field experienced by the nuclei and change the Larmor frequency as $\omega= \gamma B_0(1+\delta)$, where $\delta$ is the \emph{chemical shift} defined relative to the reference frequency $\omega_0$ to make it field independent as
\begin{equation}
	\label{eq:chemshift}
	\delta=\dfrac{\omega-\omega_0}{\omega_0}.
\end{equation}
Additionally, this energy splitting favors the spin population to be in the low energy state while the thermal energy of the system tends to equalize the population in all levels. The emerging \emph{macroscopic magnetization} $M_0$ from this population imbalance can be estimated using the \emph{Boltzmann distribution}, which describes the probability of spins occupying different energy states based on their energy and the system's temperature. At room temperature, the nuclear magnetic energy levels are much smaller than the thermal energy of the system(i.e., $\hbar\gamma B_0 \ll k_BT$, where $k_B$ is Boltzmann's constant and T is the temperature in Kelvin). With this assumption, the infinite series from Boltzmann distribution for the observable \emph{macroscopic magnetization} $M_0$ that is aligned along the $B_0$ in an \gls{NMR} experiment~\cite{haackeMagneticResonanceImaging1999} simplifies to 
\begin{equation}
	M_0 \approxeq \rho_0 \dfrac{s(s+1)\gamma^2\hbar^2}{3kT}B_0 .
	\label{eq:M0}
\end{equation}
Here, $\rho_0$ is the number of spins in a unit volume. This general expression accommodates both spin-$\tfrac{1}{2}$ hydrogen nuclei ($^1H$)and spin-1 deuterium nuclei ($^2H$) investigated in this thesis. Additionally, the above expression partially explains the sensitivity improvement of an \gls{NMR} experiment at higher field strengths ~\cite{pohmannSignaltonoiseRatioMR2016}.

Although the behavior of the spin is quantum mechanical in nature, classical models are mostly sufficient in the scope of this thesis except for topics such as polarization and quadrupolar relaxation. Therefore, the following sections rely mostly on the classical picture of assuming macroscopic magnetization created by "spinning tiny magnets" in an external magnetic field, which obeys classical rotational kinematics and differential equations~\cite{hansonQuantumMechanicsNecessary2008}. The complete quantum mechanical treatment of spins and origin of \gls{NMR} signal is covered in numerous sources~\cite{levittSpinDynamicsBasics,haackeMagneticResonanceImaging1999}. 

\subsection{Excitation and Detection}
The \emph{macroscopic magnetization} $M_0$ cannot be detected in a conventional \gls{NMR} experiment. To create a detectable signal, a time varying magnetic field $B_1$, typically in the \gls{RF} frequency range, is applied orthogonally to the $B_0$ field to create \emph{transverse magnetization} $M_{xy}$ which can then be detected using a detection coil. 
In most cases, the $B_1$ field (\gls{RF} pulse) is applied at the Larmor frequency $\omega_0$ to achieve on-resonance condition ($\omega\textsubscript{RF}=\omega_0$). For an on-resonant \gls{RF} pulse $B_1(t)= b_1(t)\cos(\omega_0 t)$, the \emph{longitudinal magnetization} $M_z$, which is initially aligned along the static magnetic field $B_0$ is tipped away from $B_0$ by the flip angle $\alpha$, which is given at any time point by
\begin{equation}
	\alpha(t)=\gamma\int_{t_0}^{t} b_1(t')\text{d}t' .
\end{equation}
In many cases, introducing a rotating frame of reference that rotates about the $B_0$ field at the instantaneous RF frequency $\omega_\text{RF}$ facilitates the investigation of the influence of the \gls{RF} pulse on the magnetization. In on-resonance condition ($\omega_\text{RF}=\omega_0$), the effect of the $B_0$ field can be ignored and the magnetization rotates about the time-independent $B_1$ field in the rotating frame. If the \gls{RF} pulse is off-resonant, an apparent magnetic field of $\Delta\omega/\gamma$ is encountered along the $B_0$ direction in the rotating frame ($\omega_\text{RF}\neq\omega_0$), where $\Delta\omega=\omega_0-\omega_\text{RF}$. During the \gls{RF} pulse, the magnetization precesses around the effective field $B_\text{eff}(t)$, which is the vector sum of  $\Delta\omega/\gamma$ and $B_1(t)$. If the frequency is modulated appropriately, the magnetization follows the $B_\text{eff}$ field to create transverse magnetization. An off-resonant \gls{RF} pulse with such a frequency modulation is called adiabatic pulse~\cite{tannusAdiabaticPulses1997}. 

The precessing \emph{transverse magnetization} $M_\text{xy}=M_x+jM_y$ created by the excitation induces a voltage in the receive coil as per Faraday's law of induction. The induced voltage is amplified, demodulated to a low frequency signal in a \emph{mixer} stage, and digitized with an \gls{ADC} for further processing.   
\subsection{Relaxation}
When the macroscopic magnetization is tipped away from the thermal equilibrium, the transverse magnetization $M_{xy}$ undergoes an exponential decay and the longitudinal magnetization $M_z$ recovers exponentially until the thermal equilibrium $M_{z,\text{eq}}$ is reached. The longitudinal relaxation is due to the exchange of energy with the environment and is characterized by the time constant $T_1$. The transverse relaxation on the other hand is due the loss of coherence as the result of spin-spin interactions and microscopic field inhomogeneties. The transverse relaxation is characterized by the time constant $T_2$. In reality, there are macroscopic field inhomogeneties that cause additional dephasing, leading to the transverse magnetization decaying even faster with a time constant $T_2^*$. This $T_2^*$ decay, however, can be recovered with a spin-echo experiment~\cite{hahnSpinEchoes1950}. The relaxation times $T_1$ and $T_2$ are inherent properties of the material at a given field strength, while $T_2^*$ depends also on the experimental conditions. 

In case of spin-1 nuclei such as deuterium($^2H$), the non-spherical charge distribution of the nuclei results in a non-zero electric quadrupolar moment. Therefore, nuclei with spin quantum number $s> \frac{1}{2}$ interact with the electric field gradients of the molecular environment resulting in a much more efficient and faster longitudinal relaxation~\cite{levittSpinDynamicsBasics}.    
\subsection{Bloch Equation}
The behavior of the macroscopic magnetization under the influence of the time varying $B_1$ field, $T_1$ and $T_2$ relaxation can be comprehensively described with a set of differential equations called Bloch equation put forth by Felix Bloch~\cite{blochNuclearInduction1946}. The matrix-vector form of the time-independent Bloch equation in the rotating frame of reference reads
\begin{equation}	
\frac{d}{dt}
\begin{pmatrix} 
	M_x' \\
	 M_y' \\
	  M_z' 
	\end{pmatrix} = \begin{pmatrix} 
	-1/T_2 & \Delta & -\gamma b_1 \\
	 -\Delta & -1/T_2 & \gamma b_1 \\ 
	 \gamma b_1 & -\gamma b_1  & -1/T_1 \end{pmatrix}
	 \begin{pmatrix} M_x' \\ M_y' \\ M_z' 
	 \end{pmatrix} 
	 + \begin{pmatrix} 0 
	 	\\ 0 \\ M_0/T_1 \end{pmatrix},
	 	\label{eq:bloch}
\end{equation}
where
	\begin{equation}
		\Delta \coloneq \gamma B_0 -\omega_\text{RF}.
	\end{equation}	
$M_x'=\operatorname{Re}(e^{j\omega_\text{RF}t}M_{xy})$, $M_y'=\operatorname{Im}(e^{j\omega_\text{RF}t}M_{xy})$ and  $M_z'=M_z$ are the magnetization components in the rotating frame, $\gamma$ is the gyromagnetic ratio, $B_0$ is the static magnetic field strength, $b_1$ is the amplitude of the demodulated \gls{RF} field, $T_1$ and $T_2$ are the longitudinal and transverse relaxation times, and $M_0=M_z(t\rightarrow \infty)$ is the equilibrium magnetization. The Bloch equation can be extended further to include effects such as diffusion~\cite{torreyBlochEquationsDiffusion1956} and chemical exchange~\cite{mcconnellReactionRatesNuclear1958}.

\subsection{Image Encoding and Reconstruction}
The generated \gls{NMR} signal explained in the previous section arises from the entire sensitive volume of the receive coil. In a standard \gls{MRI} setup, the signal is spatially encoded with additional magnetic fields along the $B_0$ direction with a linear spatial variation in three orthogonal directions referred as \emph{gradients} $\textbf{G}$. In the rotating frame, the spatial phase accrual of the signal as the result of the gradient fields is

\begin{equation}
	\label{eq:phaseaccrural}
	\Delta\phi(\textbf{r},t)=-\gamma\int_{0}^{t} \textbf{G}(t')\cdot\textbf{r} \text{d}t'=-\textbf{k}(t)\cdot\textbf{r},
\end{equation}
where 
\begin{equation}
	\label{eq:kspace}
	\textbf{k}(t)=\gamma\int_{o}^{t}\textbf{G}(t')\text{d}t'.
\end{equation}

$\textbf{k}$ is the coordinate location corresponding to a particular spatial frequency in the reciprocal space of the image called \emph{k-space}. The received signal $s(t)$ of an object with a spin density distribution $\rho(\textbf{r})$ modulated by the gradient field in the receive coil with sensitivity $W(\textbf{r})$ is given by 
\begin{equation}
	\label{eq:signal}
	s(t) \propto \int_V \rho(\textbf{r})W(\textbf{r})e^{-j\textbf{k}(t)\cdot\textbf{r}} \text{d}V.
\end{equation}

If sufficient k-space samples are acquired for a given resolution and \gls{FOV} as dictated by the \emph{Nyquist-Shannon sampling theorem}~\cite{shannonCommunicationPresenceNoise1949}, an aliasing-free MR image can be reconstructed using the \gls{DFT}. In conventional \gls{MRI}, the k-space samples are acquired on a uniformly-spaced Cartesian grid in a line-by-line fashion, which allows using the much faster realization of \gls{DFT} called \gls{FFT}~\cite{cooleyAlgorithmMachineCalculation1965}. Due to the time-consuming nature of this acquisition process, several faster techniques have been developed to acquire the entire k-space with a single excitation (single-shot)~\cite{mansfieldMultiplanarImageFormation1977} or a limited number of excitations (multi-shot).